\section{Find Distance}

Using the initial positions defined in the Problem Setup – star at $\vec{r}_\star = (0,\,0)$ and planet at $\vec{r}_p = (100,\,75)$, the initial separation $r$ is computed as:

\begin{align}
r &= \sqrt{(x_p - x_\star)^2 + (y_p - y_\star)^2} \\
  &= \sqrt{(100 - 0)^2 + (75 - 0)^2} \\
  &= \sqrt{100^2 + 75^2} \\
  &= \sqrt{10000 + 5625} \\
  &= \sqrt{15625} \\
  &= 125
\end{align}

The initial separation between the star and planet is therefore $r = 125$ (in simulation units). This value feeds directly into the gravitational acceleration calculation, the escape velocity, and the assignment of the initial orbital velocity direction, all of which follow in subsequent sections.

It is worth noting that this distance is recomputed at every timestep of the Velocity Verlet integrator, since the planet's position changes continuously. The formula used remains identical; only the current coordinates of the planet are substituted.